\cotefile{83}
\scanpage{113}
pas de l'ordre dans lequel on prend les deux \deleted{\unsure{pai|paires}} couples $(D_1,D_2)$, $(D'_1,D'_2)$. On a un rapport plus intrinsèque, qui ne change pas par permutation $D_1\leftrightarrow D_2$ ou $D'_1\leftrightarrow D'_2$
\[
\rho=\lambda+\lambda^{-1}
\]
qui est défini sur $k$ (sans localiser). On retrouve $\lambda$, au passage à $\lambda^{-1}$ près, comme solution de
\[
\lambda^2-\lambda\rho+1=0\;.
\]

\deleted{\unsure{Ce|Le} $\rho$ est \unsure{part|pour}} On aura ici
\[
\rho\ne 2\qquad\text{sur chaque fibre.}
\]
i.e.
\[
\rho_0=\rho-2\qquad\begin{aligned}&\text{est $\ne0$ sur chaque fibre}\\[-2pt]&\text{i.e. est invers.}\end{aligned}
\]

On voit \deleted{que} aisément que les classes de \unsure{\deleted{pai}res|paires|\deleted{paires}}, \deleted{dans} $(H,H')$ dans $A$, \inserted{($H,H'$ étales sur $k$, $H,H'$ en pos. gén.),} modulo \unsure{conjugaisons|conjugaison} locales correspondent aux diff $\rho_0$ inv. i.e. \unsure{si on|si on a|si} deux telles paires $(H,H')$ avec inv. $\rho_0=\rho_0(H,H')$, $(\overline H,\overline H')$ avec inv. $\overline{\rho_0}$
% archive sheet number: 112
